Facial-expression recognition (FER) is commonly formulated as predicting a category from a face image, but multiple annotators need not assign the same category. FER+ records this variation by associating FER2013 images with distributions of human responses rather than relying exclusively on the original categorical labels \cite{goodfellow2013fer,barsoum2016ferplus}. Distribution-based supervision retains the relative frequency of these responses and has been investigated as an alternative to hard-label training \cite{barsoum2016ferplus,washington2021crowdsourcing}. Agreement with a response distribution and agreement with a categorical reference are nevertheless different evaluation objectives. This distinction becomes important when a classifier is allowed to withhold predictions rather than assign a category to every image.

Selective prediction evaluates the trade-off between the proportion of inputs accepted and the disagreement or error among the accepted predictions \cite{geifman2017selective,geifman2019selectivenet}. Here, coverage denotes the fraction of images accepted, and disagreement is the mean fraction of annotations that do not match the predicted category. Two models evaluated at the same coverage may accept different images. Their selective risks can therefore differ because they make different category predictions, because they select different images, or because both change. A comparison using each model's own confidence ranking does not, by itself, distinguish these possibilities. A probability vector can also move closer to the annotation distribution without changing its most probable category, so distributional agreement and category disagreement need not improve together.

The choice of training control creates a second interpretive difficulty. Replacing a hard target with an empirical vote distribution changes both the probability assigned to the dominant category and the allocation of probability among the remaining categories. An advantage over hard-label training consequently does not isolate the contribution of the observed minority-category allocation. We examine this comparison using controls that preserve the dominant target probability while redistributing the remaining probability, together with controls that account for tied maxima. All conditions within a backbone and random seed retain the same training images, weights derived from votes for the model's output categories, architecture, initialization, sample order, augmentation, and optimization budget. The largest training-target probability describes the annotations, whereas prediction confidence is computed from the fitted model. The comparisons concern specified target representations; the softening controls do not additionally match target entropy or aggregate class proportions.

The annotation vocabulary introduces a further distinction between training and evaluation. FER+ includes contempt, unknown, and not-a-face responses in addition to the seven categories predicted by the models examined here \cite{barsoum2016ferplus}. We call categories in the model's output vocabulary supported categories. Votes for the other categories cannot be matched by a seven-category prediction. They therefore contribute to disagreement even when the model chooses the most strongly supported category available to it. Restricting the evaluation to supported votes would change the reference against which the model is scored. We instead retain the complete annotation distribution and distinguish unsupported vote mass, ambiguity among supported categories, and disagreement beyond the best available supported-category prediction.

This study brings these issues together in a matched empirical evaluation on FER2013 with FER+ annotations. The primary follow-up comparison uses validation-selected Swin-T models at 80\% coverage across three paired training seeds. ResNet-18, distributional distance, fixed-epoch checkpoints, identical-image comparisons, confidence-score sensitivity and calibration transfer retain the secondary or exploratory roles specified in Sect.~3.1. The analysis addresses three research questions:
\begin{enumerate}
\item \textbf{RQ1: Target representation.} How do observed crowd targets compare with controls that preserve the highest target probability or tied maxima, under the same training protocol?
\item \textbf{RQ2: Interpretation.} How does the comparison change when models accept identical images, and when we account for the annotation reference and categories absent from the model?
\item \textbf{RQ3: Robustness and transfer.} Which findings hold across the tested backbones, confidence scores, saved model states (checkpoints), calibration transfer and external clips?
\end{enumerate}

The results show why these comparisons need to be considered together. The selective advantage of empirical-distribution supervision over hard targets becomes smaller when the dominant target probability is preserved by the control. Although distributional agreement improves across the tested within-source configurations, the selective advantage does not persist across both backbones, and comparisons on identical accepted images depend on the selecting model. Frozen evaluation on 1,440 RAVDESS clips \cite{livingstone2018ravdess} and calibration-to-test transfer further examine how the conclusions depend on the scoring reference and achieved operating point. The contribution is a controlled account of these dependencies, using established learning architectures rather than a new recognition or rejection algorithm. Section~5 develops the implications for evaluating selective predictions.

A sender-controlled messaging prototype motivated the seven-category implementation. Related applications include text-based notifications in EmotionPush \cite{wang2016emotionpush}, facial and textual signals in EmoChat \cite{chong2019emochat}, and expression-informed emoji selection \cite{liu2019selfawareness}. Agreement with expression annotations should not be equated with a person's internal emotion, intention, or sincerity \cite{barrett2019emotions}. The prototype and its unvalidated communication role are therefore documented separately in the supplement. The empirical conclusions concern annotation-based facial-expression evaluation.
